%% ma: L12-B07-0009
%% lop: 12
%% chuong: 2
%% bai: 7
%% dang: 12.7.7
%% loai: TN4
%% muc: TH
%% dap_an: "A"
%% nguon: "LaTeX giáo viên gửi 21/10/2026 (M05, Vectơ và hệ trục tọa độ trong không gian, 107 câu), Phần I Câu 40"
%% kien_thuc: "Bài 7 (tọa độ của điểm)"
%% trang_thai: chinh-thuc
%% ghi_chu: "Giáo viên đã duyệt 21/10/2026. Đề gốc không nói gốc O và trục Oz; đã bổ sung 'gốc O là điểm xuất phát', 'trục Oz hướng thẳng đứng lên trên' cho rõ (không đổi đáp án)"
\begin{ex}
Một chiếc khinh khí cầu bay lên từ mặt đất. Thiết lập hệ trục $Oxyz$ với gốc $O$ là điểm xuất phát, mặt phẳng $(Oxy)$ là mặt đất, trục $Ox$ hướng về phía Nam, trục $Oy$ hướng về phía Đông, trục $Oz$ hướng thẳng đứng lên trên. Sau một thời gian, khinh khí cầu cách điểm xuất phát $10$ km về phía Đông, $5$ km về phía Nam và ở độ cao $400$ m. Tọa độ của khinh khí cầu (đơn vị: km) là
\choice{\True $(5;10;0{,}4)$}{$(10;5;400)$}{$(-5;10;0{,}4)$}{$(5;10;400)$}
\loigiai{Phía Nam là chiều dương $Ox$ nên $x=5$; phía Đông là chiều dương $Oy$ nên $y=10$; độ cao $400$ m $=0{,}4$ km nên $z=0{,}4$.}
\end{ex}
